\section{Conclusion}
\label{sec:conclusion}

We introduce SiliQun, an open-source quantum simulation platform that unifies silicon spin qubit technologies — including SiMOS quantum dots, phosphorus-donor systems, and gate-all-around nanowires — within a physics-grounded Lindblad density-matrix framework compatible with Gymnasium. Developed using the Physics-Grounded Iterative Refinement Stack (PGIRS) methodology, our approach systematically integrates noise-aware design principles, interface abstraction, and empirical validation throughout the simulation stack. To our knowledge, this represents the most rigorously tested open-source platform for silicon spin qubit simulations; no other open-source framework currently combines silicon-spin-specific physics with a Gymnasium-compatible RL interface. The platform is validated through two replication studies and one exploratory extension across three hardware architectures and multiple reinforcement learning algorithms.

The validation results highlight SiliQun's capabilities across several dimensions. First, the Lindblad physics engine achieves high numerical accuracy in fundamental quantum benchmarks, including Rabi oscillations and $T_1$ decay processes, with maximum deviations below $10^{-12}$ in float64 precision. Second, our replication studies establish cross-technology consistency: we successfully reproduce Daraeizadeh~et~al.'s DQL/SARSA CZ gate synthesis on SiMOS platforms ($\bar{F} = 0.9996 \pm 0.0001$, absolute deviation $\Delta F = 0.0004$ from the reference) and He~et~al.'s universal state-preparation protocol on donor systems ($\bar{F} = 0.9748 \pm 0.0098$). Applying the same DQN protocol to gate-all-around nanowires yields $\bar{F} = 0.9877 \pm 0.0021$, which we believe constitutes the first machine-learning benchmark for this emerging qubit platform.

Beyond these baseline validations, our reinforcement learning experiments reveal two complementary aspects of the system's physical realism. A lightweight PPO agent achieves near-perfect fidelity ($F = 0.9996$) on 2-qubit Bell state synthesis, confirming both the engine's learnable reward signals and correct Gymnasium implementation — a result that validates interface and reward design rather than algorithmic superiority. However, the same agent manages only $F = 0.3998$ on the 4-qubit GHZ task, demonstrating physically realistic noise scaling and genuine algorithmic challenges absent specialised enhancements.

Scalability studies conducted as part of this work expose fundamental limitations in current standard approaches. While DRL baselines reliably solve all 2-qubit target states ($F > 0.99$), performance sharply declines for 3-qubit GHZ, W, and cluster states ($\bar{F} = 0.69 \pm 0.06$). Dicke-$k_3$ states prove considerably more tractable ($\bar{F} = 0.9782 \pm 0.0137$ at $n=3$), hinting at state-dependent noise susceptibility driven by the lower entanglement entropy of this family. This scalability ceiling directly motivated the development of a companion adaptive quantum control framework (under separate submission).

Architecturally, SiliQun makes two lasting contributions. Its hardware-neutral simulation substrate enables unprecedented cross-technology comparisons through a standardised Gymnasium interface, overcoming the tool fragmentation that has historically impeded silicon spin qubit research. The PGIRS methodology further provides transferable engineering principles for extending the platform to non-silicon qubit technologies without modifying the core Lindblad solver — a flexibility we have already begun exploiting in ongoing transmon research.

Four directions emerge as priorities for future work. Comparative studies between reinforcement learning and classical optimal control methods (GRAPE, CRAB) could clarify when machine learning offers practical advantages in gate synthesis. Physical validation through randomised benchmarking would help bridge the current simulation-to-reality gap. Cross-platform generalisation to transmon and trapped-ion systems would demonstrate the full scope of PGIRS's hardware-agnostic design. Finally, extending scalability studies beyond four qubits may reveal whether the observed noise-conditional hyperparameter relationships generalise to larger registers.

SiliQun is available under the MIT licence at \url{https://github.com/r3d-phd/siliqun}, with full reproducibility ensured through provided scripts and archived Aziz HPC job files. We encourage community contributions to expand both its physical models and its pedagogical applications.
